% ==================== 全局可调参数 ====================
% 修改方法：优先改这里，能统一控制整份简历的字号和行距。
% 例如：11pt -> 10pt / 12pt；\ResumeLineSpacing{} 调大可更松，调小可更紧凑。
\documentclass[11pt,a4paper]{article}

% 引入中文支持包
\usepackage[UTF8]{ctex}
% 页面边距设置 (一页纸简历的核心，充分利用空间)
\usepackage[left=1.5cm, right=1.5cm, top=1.5cm, bottom=1.5cm]{geometry}
% 标题格式自定义
\usepackage{titlesec}
% 列表格式自定义
\usepackage{enumitem}
% 行距控制
\usepackage{setspace}
% 线性图标（接近示例风格）
\usepackage{fontawesome5}
% 图片支持
\usepackage{graphicx}
% 图形裁切（用于圆形校徽）
\usepackage{tikz}
% 超链接支持
\usepackage[colorlinks=true, linkcolor=blue, urlcolor=black]{hyperref}
% 颜色包
\usepackage{xcolor}

% --- 样式自定义 ---
% 定义主题色 (深蓝色，显得专业沉稳)
\definecolor{ThemeColor}{RGB}{25, 75, 125}
\definecolor{LightGray}{RGB}{225, 231, 238}

% ==================== 字体大小 / 行距 ====================
% 修改方法：
% 1. 整体字号：改 documentclass 里的 11pt，例如 10pt 更紧凑，12pt 更大。
% 2. 局部字号：把 \small / \normalsize / \large / \LARGE 改成你想要的大小。
% 3. 行距：这里的 1.0 = 标准行距，1.05 ~ 1.15 更适合简历。
\newcommand{\ResumeLineSpacing}{1.06}
\setstretch{\ResumeLineSpacing}

% 需要更细粒度控制时，可以直接在某一段前使用：
% \begin{spacing}{1.08}
%   这一段会单独变松
% \end{spacing}

% 自定义 section 标题样式：大号字体、加粗、主题色、带下划线
\titleformat{\section}{\large\bfseries\color{ThemeColor}}{}{0em}{}[\vspace{-1pt}\color{LightGray}\titlerule]
\titlespacing{\section}{0pt}{10pt}{5pt} % 控制标题上下的间距

% 标题字号修改方法：把 \large 改成 \normalsize / \small / \Large / \LARGE
% 如果标题之间太挤，可以优先调 \titlespacing 的上下间距

% 自定义 itemize 列表样式：取消多余缩进，减小行距
% 修改方法：
% - noitemsep：去掉条目之间的额外空隙
% - topsep：列表上下留白
% - parsep：同一条目内多段落的间距
\setlist[itemize]{leftmargin=*, noitemsep, topsep=2pt, partopsep=0pt, parsep=2pt}

% 标题前的小图案（统一风格）
\newcommand{\titleicon}[1]{\textcolor{ThemeColor}{\normalsize #1}\hspace{0.35em}}
\newcommand{\resumesection}[2]{\section*{\titleicon{#1}#2}}
\newcommand{\eduicon}{\faIcon{graduation-cap}}
\newcommand{\skillicon}{\faIcon{wrench}}
\newcommand{\projicon}{\faIcon{code}}
\newcommand{\awardicon}{\faIcon{trophy}}
\newcommand{\selficon}{\faIcon{user}}
% 新增可选区块图标
\newcommand{\internicon}{\faIcon{briefcase}}
\newcommand{\campusicon}{\faIcon{users}}
\newcommand{\langicon}{\faIcon{language}}
\newcommand{\researchicon}{\faIcon{microscope}}
\newcommand{\topicicon}{\faIcon{book}}

% 图片占位或加载（照片/校徽）
\newcommand{\imageorplaceholder}[3]{
    \IfFileExists{#1}{
        \includegraphics[width=#2,height=#3,keepaspectratio]{#1}
    }{
        \fcolorbox{LightGray}{white}{\parbox[c][#3][c]{#2}{\centering\scriptsize\color{gray}图片占位}}
    }
}

% 圆形图标（用于学校图标）
\newcommand{\circleimageorplaceholder}[2]{
    \IfFileExists{#1}{
        \begin{tikzpicture}
            \clip (0,0) circle (#2/2);
            \node at (0,0) {\includegraphics[width=#2,height=#2]{#1}};
            \draw[LightGray, line width=0.3pt] (0,0) circle (#2/2);
        \end{tikzpicture}
    }{
        \begin{tikzpicture}
            \draw[LightGray, line width=0.4pt, fill=white] (0,0) circle (#2/2);
            \node[gray] at (0,0) {\scriptsize 校徽};
        \end{tikzpicture}
    }
}

\begin{document}
% 取消页码
\pagestyle{empty}

% 如果你想让正文整体再紧凑一点，可在这里改成：
% \small
% 如果想整体更大一点，可改成：
% \normalsize

% ==================== 头部信息 ====================
\noindent
\begin{minipage}[c]{0.17\textwidth}
    ~
\end{minipage}%
\hfill
\begin{minipage}[c]{0.66\textwidth}
    \centering
    % 名字字号修改：\LARGE / \Large / \huge 都可以试，越往后越大
    {\LARGE \textbf{你的名字}} \\[6pt]
    % 这里控制头部个人信息的字号，改成 \normalsize 会更大，改成 \footnotesize 会更紧凑
    \small
    \textbf{求职意向：} C/C++ 后台开发日常实习生 \\[3pt]
    \textbf{到岗时间：} 每周可实习4-5天，可随时入职 \\[4pt]
    \textcolor{ThemeColor}{\faIcon{map-marker-alt}}\ 深圳南山区 \quad|\quad
    \textcolor{ThemeColor}{\faIcon{phone}}\ 138-XXXX-XXXX \\[2pt]
    \textcolor{ThemeColor}{\faIcon{at}}\ your.email@example.com \quad|\quad
    \textcolor{ThemeColor}{\faIcon{github}}\ \href{https://github.com/your_id}{github.com/your\_id}
\end{minipage}%
\hfill
\begin{minipage}[c]{0.15\textwidth}
    \raggedleft\imageorplaceholder{assets/images/photo.jpg}{2.0cm}{2.55cm}
\end{minipage}

\vspace{5pt}

% ==================== 教育背景 ====================
\resumesection{\eduicon}{教育背景}
\noindent\textbf{某某大学\quad 计算机科学与技术\quad 本科 (大二在读)} \hfill 2024年9月 -- 2028年6月
\begin{itemize}
    \item \textbf{主修课程：} 数据结构、算法设计与分析、C++面向对象编程等。
    \item \textbf{学业进展：} 正在系统深入学习 Linux 网络编程、操作系统底层原理及现代 C++ 特性。
\end{itemize}

% ==================== 专业技能 ====================
\resumesection{\skillicon}{专业技能}
\begin{itemize}
    \item \textbf{编程语言：} 熟悉 C/C++，了解面向对象编程思想；熟悉常用 STL 容器及其底层原理；正在学习 C++11 新特性（智能指针、右值引用等）与内存管理。
    \item \textbf{数据结构与算法：} 具备扎实的算法功底，熟练掌握排序、贪心、二分、枚举、树等算法；有蓝桥杯省赛备赛经验，目前持续在 LeetCode 进行针对性刷题训练。
    \item \textbf{计算机基础：} 了解计算机网络（TCP/IP模型、TCP三次握手等）与操作系统（进程与线程概念）的基础理论，目前正在结合项目深入学习。
    \item \textbf{开发环境：} 熟悉 Linux 基础环境配置与常用命令，了解基本的 Git 版本控制操作。
\end{itemize}

% ==================== 项目经历 ====================
\resumesection{\projicon}{项目经历}
\textbf{轻量级高性能 Web 服务器 (开发中)} \hfill 个人独立项目 \\
\textit{C++, Linux, epoll, Socket} \hfill 2026年5月 -- 至今
\begin{itemize}
    \item \textbf{项目背景：} 为加深对网络编程和操作系统底层机制的理解，基于 Linux 环境和 C++ 开发的高并发 Web 服务器。
    \item \textbf{核心技术：} 采用 \texttt{epoll} I/O 多路复用技术处理高并发客户端连接；实现基于有限状态机的 HTTP 请求解析。
    \item \textbf{当前进展：} 已完成基础的 TCP Socket 通信模块封装，正在接入 \texttt{epoll} 并发模型与线程池，后续将引入压测工具进行性能调优。
\end{itemize}

% ==================== 荣誉奖项 ====================
\resumesection{\awardicon}{荣誉奖项}
\begin{itemize}
    \item 第XX届 蓝桥杯全国软件和信息技术专业人才大赛 (C/C++大学组) \textbf{省赛X等奖} \hfill 2025年X月
    \item 某某大学 优秀学生/校级奖学金 (若有则写，无则删除此行) \hfill 2025年X月
\end{itemize}

% ==================== 自我评价 ====================
\resumesection{\selficon}{自我评价}
\begin{itemize}
    \item \textbf{极强的学习自驱力：} 明确大厂 C/C++ 岗位要求，大二起主动脱离“舒适区”，制定了高强度的底层技术补齐计划，能够快速吸收并应用新技术。
    \item \textbf{地理与时间优势：} 身处深圳南山，时间充裕，无入职缓冲期。对技术充满热情，极其渴望进入真实的工业级团队打磨工程能力，能承受较高强度的工作节奏。
\end{itemize}

% --------------------------------------------------
% 可选区块模板（注释掉的示例）
% 使用方法：将下面对应区块前面的百分号（%）去掉，复制到需要的位置，填写具体内容。
% --------------------------------------------------

% % \resumesection{\internicon}{实习经历}
% % \textbf{公司名称} \hfill 城市 \\ 
% % 职位：实习生 \quad 年份：2023.07 -- 2024.04
% % \begin{itemize}
% %     \item 主要职责：例如负责后端模块开发、需求拆解与实现。
% %     \item 成果/量化：例如优化某接口，吞吐提升 X%，或修复 Y 个关键 bug。
% % \end{itemize}
%
% % \resumesection{\campusicon}{校园经历}
% % \textbf{学生组织 / 社团名称} \hfill 学校 \\ 
% % 职务：负责人/成员 \quad 年份：2022 -- 2023
% % \begin{itemize}
% %     \item 活动或项目：例如组织技术沙龙、带领团队完成比赛。
% %     \item 贡献与能力：例如项目管理、协调资源、活动策划。
% % \end{itemize}
%
% % \resumesection{\langicon}{语言技能}
% % \begin{itemize}
% %     \item 汉语：母语
% %     \item 英语：阅读（专业），听说（良好）
% %     \item 其他语言：如日语/法语等，写明水平等级或证书。
% % \end{itemize}
%
% % \resumesection{\researchicon}{科研成果}
% % \begin{itemize}
% %     \item 论文/成果名称（期刊/会议，年份）—— 简短说明贡献与成果。
% %     \item 项目或成果：如算法改进、数据集发布、工具/库的开源等。
% % \end{itemize}
%
% % \resumesection{\topicicon}{研究课题}
% % \begin{itemize}
% %     \item 课题名称：简述研究目标和方法。
% %     \item 目前进展：例如已完成的工作（数据收集、模型搭建、实验验证）。
% %     \item 未来计划：后续要做的实验或预期发表计划。
% % \end{itemize}

\end{document}